\documentclass[UTF8,12pt]{ctexart}
\usepackage[a4paper,margin=24mm]{geometry}
\usepackage{amsmath,amssymb,bm,graphicx,adjustbox}
\newcommand{\fitmath}[1]{\begin{adjustbox}{max width=\linewidth}$\displaystyle #1$\end{adjustbox}}
\setlength{\parindent}{0pt}
\setlength{\parskip}{.7em}
\sloppy
\allowdisplaybreaks
\begin{document}
\section*{2001 年考研数学三 · 真题}
\subsection*{2001年全国硕士研究生招生考试试题}

\subsection*{一、填空题（本题共5小题，每小题3分，满分15分）}

（1）设生产函数为 \(\displaystyle Q=AL^\alpha K^\beta\)，其中 \(\displaystyle Q\) 是产出量，\(\displaystyle L\) 是劳动投入量，\(\displaystyle K\) 是资本投入量，而 \(\displaystyle A,\allowbreak \alpha,\allowbreak \beta\) 均为大于零的参数，则当 \(\displaystyle Q=1\) 时 \(\displaystyle K\) 关于 \(\displaystyle L\) 的弹性为\_\_\_\_\_\_。


（2）某公司每年的工资总额在比上一年增加20\%的基础上再追加2百万元。若以 \(\displaystyle W_t\) 表示第 \(\displaystyle t\) 年的工资总额（单位：百万元），则 \(\displaystyle W_t\) 满足的差分方程是\_\_\_\_\_\_。


（3）设矩阵 \(\displaystyle A=\begin{pmatrix}k&1&1&1\\1&k&1&1\\1&1&k&1\\1&1&1&k\end{pmatrix}\)，且 \(\displaystyle r(A)=3\)，则 \(\displaystyle k=\)\_\_\_\_\_\_。


（4）设随机变量 \(\displaystyle X\) 和 \(\displaystyle Y\) 的数学期望分别为 \(\displaystyle -2\) 和 \(\displaystyle 2\)，方差分别为 \(\displaystyle 1\) 和 \(\displaystyle 4\)，而相关系数为 \(\displaystyle -0.5\)，则根据切比雪夫不等式 \(\displaystyle P\{|X+Y|\ge6\}\le\)\_\_\_\_\_\_。


（5）设总体 \(\displaystyle X\) 服从正态分布 \(\displaystyle N(0,\allowbreak 2^2)\)，而 \(\displaystyle X_1,\allowbreak X_2,\allowbreak \ldots,\allowbreak X_{15}\) 是来自总体 \(\displaystyle X\) 的简单随机样本，则随机变量 \(\displaystyle Y=\dfrac{\displaystyle X_1^2+\cdots+X_{10}^2}{\displaystyle 2(X_{11}^2+\cdots+X_{15}^2)}\) 服从\_\_\_\_\_\_分布，参数为\_\_\_\_\_\_。


\subsection*{二、选择题（本题共5小题，每小题3分，满分15分）}

（1）设 \(\displaystyle f(x)\) 的导数在 \(\displaystyle x=a\) 处连续，又 \(\displaystyle \lim_{x\to a}\dfrac{\displaystyle f^{\prime}(x)}{\displaystyle x-a}=-1\)，则（　）。
（A）\(\displaystyle x=a\) 是 \(\displaystyle f(x)\) 的极小值点。（B）\(\displaystyle x=a\) 是 \(\displaystyle f(x)\) 的极大值点。（C）\(\displaystyle (a,\allowbreak f(a))\) 是曲线 \(\displaystyle y=f(x)\) 的拐点。（D）\(\displaystyle x=a\) 不是 \(\displaystyle f(x)\) 的极值点，\(\displaystyle (a,\allowbreak f(a))\) 也不是曲线 \(\displaystyle y=f(x)\) 的拐点。


（2）设 \(\displaystyle g(x)=\displaystyle \int_0^x f(u)\,\mathrm du\)，其中 \(\displaystyle f(x)=\begin{cases}\dfrac{\displaystyle 1}{\displaystyle 2}(x^2+1),\allowbreak &0\le x<1,\allowbreak \\\dfrac{\displaystyle 1}{\displaystyle 3}(x-1),\allowbreak &1\le x\le2,\allowbreak \end{cases}\) 则 \(\displaystyle g(x)\) 在区间 \(\displaystyle (0,\allowbreak 2)\) 内（　）。
（A）无界。（B）递减。（C）不连续。（D）连续。


（3）设 \(\displaystyle A=\begin{pmatrix}a_{11}&a_{12}&a_{13}&a_{14}\\a_{21}&a_{22}&a_{23}&a_{24}\\a_{31}&a_{32}&a_{33}&a_{34}\\a_{41}&a_{42}&a_{43}&a_{44}\end{pmatrix}\)，\(\displaystyle B=\begin{pmatrix}a_{14}&a_{13}&a_{12}&a_{11}\\a_{24}&a_{23}&a_{22}&a_{21}\\a_{34}&a_{33}&a_{32}&a_{31}\\a_{44}&a_{43}&a_{42}&a_{41}\end{pmatrix}\)，\(\displaystyle P_1=\begin{pmatrix}0&0&0&1\\0&1&0&0\\0&0&1&0\\1&0&0&0\end{pmatrix}\)，\(\displaystyle P_2=\begin{pmatrix}1&0&0&0\\0&0&1&0\\0&1&0&0\\0&0&0&1\end{pmatrix}\)，其中 \(\displaystyle A\) 可逆，则 \(\displaystyle B^{-1}\) 等于（　）。


（3）（续）（A）\(\displaystyle A^{-1}P_1P_2\)。（B）\(\displaystyle P_1A^{-1}P_2\)。（C）\(\displaystyle P_1P_2A^{-1}\)。（D）\(\displaystyle P_2A^{-1}P_1\)。


（4）设 \(\displaystyle A\) 是 \(\displaystyle n\) 阶矩阵，\(\displaystyle \alpha\) 是 \(\displaystyle n\) 维列向量。若 \(\displaystyle r\!\begin{pmatrix}A&\alpha\\\alpha^{\mathrm T}&0\end{pmatrix}=r(A)\)，则线性方程组（　）。
（A）\(\displaystyle AX=\alpha\) 必有无穷多解。（B）\(\displaystyle AX=\alpha\) 必有唯一解。（C）\(\displaystyle \begin{pmatrix}A&\alpha\\\alpha^{\mathrm T}&0\end{pmatrix}\begin{pmatrix}X\\y\end{pmatrix}=0\) 仅有零解。（D）\(\displaystyle \begin{pmatrix}A&\alpha\\\alpha^{\mathrm T}&0\end{pmatrix}\begin{pmatrix}X\\y\end{pmatrix}=0\) 必有非零解。


（5）将一枚硬币重复掷 \(\displaystyle n\) 次，以 \(\displaystyle X\) 和 \(\displaystyle Y\) 分别表示正面向上和反面向上的次数，则 \(\displaystyle X\) 和 \(\displaystyle Y\) 的相关系数等于（　）。
（A）\(\displaystyle -1\)。（B）\(\displaystyle 0\)。（C）\(\displaystyle \dfrac{\displaystyle 1}{\displaystyle 2}\)。（D）\(\displaystyle 1\)。


\subsection*{三、（本题满分5分）}

设 \(\displaystyle u=f(x,\allowbreak y,\allowbreak z)\) 有连续的一阶偏导数，又函数 \(\displaystyle y=y(x)\) 及 \(\displaystyle z=z(x)\) 分别由下列两式确定：\(\displaystyle e^{xy}-xy=2\) 和 \(\displaystyle e^x=\displaystyle \int_0^{x-z}\dfrac{\displaystyle \sin t}{\displaystyle t}\,\mathrm dt\)，求 \(\displaystyle \dfrac{\displaystyle \mathrm du}{\displaystyle \mathrm dx}\)。


\subsection*{四、（本题满分6分）}

已知 \(\displaystyle f(x)\) 在 \(\displaystyle (-\infty,\allowbreak +\infty)\) 内可导，且 \(\displaystyle \lim_{x\to\infty}f^{\prime}(x)=e\)，\(\displaystyle \lim_{x\to\infty}\left(\dfrac{\displaystyle x+c}{\displaystyle x-c}\right)^x=\displaystyle \lim_{x\to\infty}[f(x)-f(x-1)]\)，求 \(\displaystyle c\) 的值。


\subsection*{五、（本题满分6分）}

求二重积分 \(\displaystyle \iint_D y[1+xe^{\dfrac{\displaystyle 1}{\displaystyle 2}(x^2+y^2)}]\,\mathrm dx\mathrm dy\) 的值，其中 \(\displaystyle D\) 是由直线 \(\displaystyle y=x\)、\(\displaystyle y=-1\) 及 \(\displaystyle x=1\) 围成的平面区域。


\subsection*{六、（本题满分7分）}

已知抛物线 \(\displaystyle y=px^2+qx\)（其中 \(\displaystyle p<0,\allowbreak q>0\)）在第一象限内与直线 \(\displaystyle x+y=5\) 相切，且此抛物线与 \(\displaystyle x\) 轴所围成的平面图形的面积为 \(\displaystyle S\)。（1）问 \(\displaystyle p\) 和 \(\displaystyle q\) 为何值时，\(\displaystyle S\) 达到最大值？（2）求出此最大值。


\subsection*{七、（本题满分6分）}

设 \(\displaystyle f(x)\) 在 \(\displaystyle [0,\allowbreak 1]\) 上连续，在 \(\displaystyle (0,\allowbreak 1)\) 内可导，且满足 \(\displaystyle f(1)=k\displaystyle \int_0^{1/k}xe^{1-x}f(x)\,\mathrm dx\;(k>1)\)，证明至少存在一点 \(\displaystyle \xi\in(0,\allowbreak 1)\)，使得 \(\displaystyle f^{\prime}(\xi)=(1-\xi^{-1})f(\xi)\)。


\subsection*{八、（本题满分7分）}

已知 \(\displaystyle f_n(x)\) 满足 \(\displaystyle f_n^{\prime}(x)=f_n(x)+x^{n-1}e^x\)（\(\displaystyle n\) 为正整数），且 \(\displaystyle f_n(1)=\dfrac{\displaystyle e}{\displaystyle n}\)，求函数项级数 \(\displaystyle \sum_{n=1}^{\infty}f_n(x)\) 之和。


\subsection*{九、（本题满分9分）}

设矩阵 \(\displaystyle A=\begin{pmatrix}1&1&a\\1&a&1\\a&1&1\end{pmatrix}\)，\(\displaystyle \beta=\begin{pmatrix}1\\1\\-2\end{pmatrix}\)。已知线性方程组 \(\displaystyle AX=\beta\) 有解但不唯一，试求（1）\(\displaystyle a\) 的值；（2）正交矩阵 \(\displaystyle Q\)，使 \(\displaystyle Q^{\mathrm T}AQ\) 为对角矩阵。


\subsection*{十、（本题满分8分）}

设 \(\displaystyle A\) 为 \(\displaystyle n\) 阶实对称矩阵，\(\displaystyle r(A)=n\)，\(\displaystyle A_{ij}\) 是 \(\displaystyle A=(a_{ij})_{n\times n}\) 中元素 \(\displaystyle a_{ij}\) 的代数余子式（\(\displaystyle i,\allowbreak j=1,\allowbreak 2,\allowbreak \cdots,\allowbreak n\)），二次型 \(\displaystyle f(x_1,\allowbreak x_2,\allowbreak \cdots,\allowbreak x_n)=\displaystyle \sum_{i=1}^n\displaystyle \sum_{j=1}^n\dfrac{\displaystyle A_{ij}}{\displaystyle |A|}x_ix_j\)。（1）记 \(\displaystyle X=(x_1,\allowbreak x_2,\allowbreak \cdots,\allowbreak x_n)^{\mathrm T}\)，把 \(\displaystyle f(x_1,\allowbreak x_2,\allowbreak \cdots,\allowbreak x_n)\) 写成矩阵形式，并证明二次型 \(\displaystyle f(X)\) 的矩阵为 \(\displaystyle A^{-1}\)；（2）二次型 \(\displaystyle g(X)=X^{\mathrm T}AX\) 与 \(\displaystyle f(X)\) 的规范形是否相同？说明理由。


\subsection*{十一、（本题满分8分）}

一生产线生产的产品成箱包装，每箱的重量是随机的。假设每箱平均重50千克，标准差为5千克。若用最大载重量为5吨的汽车承运，试利用中心极限定理说明每辆车最多可以装多少箱，才能保障不超载的概率大于0.977。（\(\displaystyle \Phi(2)=0.977\)，其中 \(\displaystyle \Phi(x)\) 是标准正态分布函数。）


\subsection*{十二、（本题满分8分）}

设随机变量 \(\displaystyle X\) 和 \(\displaystyle Y\) 的联合分布是正方形 \(\displaystyle G=\{(x,\allowbreak y)\mid1\le x\le3,\allowbreak 1\le y\le3\}\) 上的均匀分布，试求随机变量 \(\displaystyle U=|X-Y|\) 的概率密度 \(\displaystyle p(u)\)。

\end{document}
